%% ===================================================================
%%  Comment on "Multiobjective model predictive control"
%%  [Automatica 45 (2009) 2823-2830]
%%  ------------------------------------------------------------------
%%  Compile:  latexmk -pdf comment.tex
%% ===================================================================
\documentclass[review]{elsarticle}

\usepackage{amsmath,amssymb,amsthm}
\usepackage[colorlinks=true,linkcolor=black,citecolor=black,urlcolor=blue]{hyperref}

\newtheorem{lemma}{Lemma}
\newproof{proof}{Proof}

\newcommand{\Rset}{\mathbb{R}}

\journal{Automatica}

\begin{document}
\begin{frontmatter}

\title{Comment on ``Multiobjective model predictive control''
[Automatica 45 (2009) 2823--2830]}

\author[ual]{Francisco Manuel Arrabal-Campos\corref{cor}}
\ead{fmarrabal@ual.es}
\author[ual]{V\'ictor Valdivieso}
\author[ual]{Emilio L\'opez-Lao}
\author[cuc]{Alejandro Campa-Pinto}
\address[ual]{Department of Engineering, Research Centre CIAIMBITAL,
University of Almer\'ia, 04120 Almer\'ia, Spain}
\address[cuc]{Department of Computer Science and Electronics,
Universidad de la Costa, Barranquilla 080002, Colombia}
\cortext[cor]{Corresponding author.}

\begin{abstract}
Lemma~4 of the commented paper states that the value function of the
multiparametric linear program with weight vector $\mu$ in the cost and state
$x$ in the right-hand side is convex and piecewise affine in $\mu$ for fixed
$x$. The function is piecewise affine, but it is \emph{concave} in $\mu$, not
convex: for each fixed decision variable the objective is affine in $\mu$ and
the feasible set does not depend on $\mu$, so the value function is a pointwise
minimum of affine functions. We trace the slip to a duality remark in the
companion conference paper, and we work out the consequences. The
weight-selection linear program of Theorem~6 evaluates the value function as
the \emph{maximum} of its affine pieces, so it describes a convex inner
approximation of the true admissible weight set, which is in general
non-convex: every weight it returns satisfies the contractive constraint, so
all closed-loop stability guarantees of the commented paper are preserved, but
the claimed equivalence does not hold, and the program can be infeasible while
admissible weights exist. In the quadratic case, the joint convexity invoked in
Theorem~7 from a 1964 result of Mangasarian and Rosen does not apply, since
that result concerns right-hand-side perturbations; the online problem is not a
convex program. The practical remedy is to treat the admissible weight set by
enumeration or gridding rather than by convex programming.
\end{abstract}

\begin{keyword}
Predictive control \sep multiobjective optimization \sep multiparametric
programming
\end{keyword}

\end{frontmatter}

% ====================================================================
\section{The claim}
% ====================================================================
The commented paper \citep{bemporad2009} builds a multiobjective MPC scheme in
which, at each sample, a weight vector is selected online subject to the
contractive constraint $V^\ast(\mu,x)\le J_a$, where $V^\ast$ is the optimal
value of a multiparametric program with the weights in the cost. For the
piecewise affine case, the machinery rests on
\citep[Lemma~4]{bemporad2009}, which considers the multiparametric linear
problem---their problem~(13)---with parameters $\mu\in\Rset^l$ in the cost
function and $x\in\Rset^n$ in the right-hand side of the constraints, and
states that the value function $V^\ast$ is
\begin{quote}
``continuous w.r.t.\ $(\mu,x)$, convex and piecewise affine w.r.t.\ $\mu$ for
any given $x$ and w.r.t.\ $x$ for any given $\mu$.''
\end{quote}
The statement for $x$ is correct, and so is piecewise affinity in $\mu$. The
convexity claim for $\mu$, however, has the wrong sign.

% ====================================================================
\section{The correction}
% ====================================================================
\begin{lemma}\label{lem:concave}
Let $V^\ast(\mu,x)=\min_z\{(c_0+C_\mu'\mu)'z:\ Gz\le b+Sx\}$ be the value
function of problem (13) of \citep{bemporad2009}. For each fixed $x$ such that
the feasible set is nonempty and the minimum is attained,
$\mu\mapsto V^\ast(\mu,x)$ is concave and piecewise affine on its domain.
\end{lemma}

\begin{proof}
For each fixed feasible $z$ the map $\mu\mapsto(c_0+C_\mu'\mu)'z$ is affine.
The feasible set $\{z:Gz\le b+Sx\}$ does not depend on $\mu$. Hence
$V^\ast(\cdot,x)$ is a pointwise minimum of a family of affine functions of
$\mu$, and therefore concave. Since for a linear program the minimum is
attained at one of finitely many vertices (or the problem is unbounded),
$V^\ast(\cdot,x)$ is the minimum of finitely many affine functions, hence
piecewise affine.
\end{proof}

This is the standard sensitivity dichotomy of parametric optimization: under
minimization the optimal value is convex in right-hand-side perturbations and
concave in parameters that enter the objective linearly
\citep[Sec.~5.6.1, Ex.~3.9]{boyd2004}. Both signs invert under maximization.
The origin of the slip appears to be the companion conference paper
\citep[p.~2406]{bemporad2009ecc}, which establishes the properties in $x$ and
adds that ``by duality, the same properties hold with respect to $\mu$.''
Duality does exchange right-hand-side and cost parameters, but it also turns
the primal minimum into a dual maximum, and that is precisely the step that
flips the curvature.

We also verified the sign numerically on quadratic instances with the same
parametric structure (weights entering the cost linearly, feasible set
independent of the weights): over $300$ randomly generated chord tests the
concavity inequality held in all cases and the convexity inequality failed in
all cases, while the same test applied to $x$ confirmed convexity, exactly as
in the corrected statement.

% ====================================================================
\section{Consequences within the commented paper}
% ====================================================================

\paragraph{Theorem 6 and the weight-selection LP (17)}
The proof of Theorem~6 in \citep{bemporad2009} invokes the result of
\citet{schechter1987} for \emph{convex} piecewise affine functions to evaluate
$V^\ast(\mu,x)$, for fixed $x$, as the \emph{maximum} of the affine pieces
$\{(\phi_i x+\gamma_i)'(c_0'+C_\mu'\mu)\}_{i\in I(x)}$, and the linear program
(17) accordingly imposes
\[
(\phi_i x+\gamma_i)'(c_0'+C_\mu'\mu)\ \le\ J_a
\qquad\text{for all } i\in I(x).
\]
By Lemma~\ref{lem:concave} the value function is the \emph{minimum} of those
pieces. Requiring every piece to lie below $J_a$ is therefore sufficient but
not necessary for $V^\ast(\mu,x)\le J_a$: the constraint set of (17) is a
convex \emph{inner approximation} of the true admissible weight set
\[
\mathcal{A}(x,J_a)=\{\mu:\ V^\ast(\mu,x)\le J_a\},
\]
which, being a sublevel set of a concave function, is in general non-convex
(for scalar $\mu$ it is a union of at most two intervals whenever the maximum
of $V^\ast(\cdot,x)$ is interior and exceeds $J_a$).

Two things follow, one reassuring and one not. Since every weight accepted by
(17) satisfies the contractive constraint, all closed-loop stability
guarantees built on that constraint are untouched; the scheme is safe as
published. What fails is exactness: (17) is not equivalent to the weight
selection problem it is meant to solve, it can be infeasible while
$\mathcal{A}(x,J_a)$ is nonempty, and its optimizer need not be the closest
admissible weight to the desired one. In our numerical experiments with the
corrected sign, the admissible set was non-convex for roughly one fifth of the
contraction levels sampled, so the gap is not a measure-zero curiosity.

\paragraph{Theorem 7 and the quadratic case}
For the branch with one quadratic objective, the proof of Theorem~7 in
\citep{bemporad2009} asserts joint convexity of $V^\ast_\mu(\mu,x)$ in
$(\mu,x)$ ``by the results of Mangasarian and Rosen (1964),'' and concludes
that the online problem (21) is a convex program. In problem (20), however,
$\mu$ enters the objective linearly through the term $\mu'C_\mu z$ while the
constraints $Gz\le b+Sx$ do not depend on $\mu$; the situation covered by
\citet{mangasarian1964} is perturbation of the constraint right-hand side, for
which the value function of a minimization problem is indeed convex. In the
$\mu$ direction the same argument as in Lemma~\ref{lem:concave} applies:
$V^\ast_\mu(\cdot,x)$ is a pointwise minimum of functions affine in $\mu$,
hence concave, and joint convexity in $(\mu,x)$ cannot hold except in
degenerate cases. Consequently the feasible set of (21) is in general
non-convex and (21) is not a convex program. Unlike the piecewise affine case,
no inner-approximation safety net is created here; any candidate weight
returned by a solver should therefore be certified a posteriori by evaluating
$V^\ast_\mu$ directly, which is cheap.

\paragraph{Practical remedy}
The admissible weight set can be handled by enumeration over the finitely many
affine pieces, or by gridding the simplex and testing
$V^\ast(\mu,x)\le J_a$ pointwise; both preserve exactness at modest cost for
the low weight dimensions typical of multiobjective MPC. Projection or
gradient steps on the weight, and any construction that presumes convexity of
$\mathcal{A}(x,J_a)$, can silently leave the admissible set.

% ====================================================================
\section*{Closing remark}
% ====================================================================
The contribution of the commented paper---a multiobjective MPC formulation
with closed-loop guarantees obtained through a contractive constraint on a
weighted value function---is not in question, and the published algorithm
remains safe precisely because the affected construction errs on the
conservative side. The correction concerns the exactness claims attached to
the weight-selection subproblems. The authors were contacted about this
observation prior to submission.

\bibliographystyle{elsarticle-harv}
\begin{thebibliography}{5}

\bibitem[Bemporad and Mu\~noz de la Pe\~na(2009a)]{bemporad2009}
Bemporad, A., Mu\~noz de la Pe\~na, D., 2009a.
Multiobjective model predictive control.
Automatica 45~(12), 2823--2830.

\bibitem[Bemporad and Mu\~noz de la Pe\~na(2009b)]{bemporad2009ecc}
Bemporad, A., Mu\~noz de la Pe\~na, D., 2009b.
Multiobjective model predictive control based on convex piecewise affine costs.
In: Proceedings of the European Control Conference (ECC). Budapest, Hungary,
pp. 2402--2407.

\bibitem[Boyd and Vandenberghe(2004)]{boyd2004}
Boyd, S., Vandenberghe, L., 2004.
Convex Optimization.
Cambridge University Press, Cambridge, UK.

\bibitem[Mangasarian and Rosen(1964)]{mangasarian1964}
Mangasarian, O.~L., Rosen, J.~B., 1964.
Inequalities for stochastic nonlinear programming problems.
Operations Research 12~(1), 143--154.

\bibitem[Schechter(1987)]{schechter1987}
Schechter, M., 1987.
Polyhedral functions and multiparametric linear programming.
Journal of Optimization Theory and Applications 53~(2), 269--280.

\end{thebibliography}

\end{document}
